\chapter{Project 2: Red Light, Green Light}

\section{Overview}
In this project, you will build a reflex game with a button and a 3-pixel WS2812B LED strip.
Each reaction window cycles through red, yellow, and green phases:
\begin{itemize}
    \item During \textbf{RED}, do not press. Pressing during red ends the round immediately.
    \item During \textbf{YELLOW}, get ready for green and do not press yet.
    \item During \textbf{GREEN}, press the button as quickly as possible to score a point.
\end{itemize}

Each successful green press prints your reaction time in milliseconds to the serial console.
Each round has 5 total reaction windows. A hit during green scores a point, and a missed green
window counts as a miss.

\begin{tcolorbox}[colback=blue!5!white,colframe=blue!40!black]
    \textbf{Learning goals:}
    \begin{itemize}
        \item Read a pushbutton using an internal pull-up resistor
        \item Control WS2812B addressable LEDs with MicroPython
        \item Build a simple game state machine
        \item Measure reaction time with millisecond timers
    \end{itemize}
\end{tcolorbox}

\pagebreak

\section{Components}
\begin{itemize}
    \item Seeed XIAO ESP32-C3 microcontroller
    \item 3x WS2812B addressable RGB LEDs wired in series
    \item 1x pushbutton
    \item Breadboard
    \item Jumper wires
\end{itemize}

\section{Directions}

\subsubsection{Remove previous components}
Before starting Project 2, remove components and wires from the previous project.

\subsection{Creating the circuit}

\subsubsection{WS2812B strip wiring}
The strip has a right-angle header soldered to its \textbf{DI} end, so it stands upright once
seated. Push those three pins into the top half of the board: \textbf{5V} into
\bbhole{pixels.header.power}{J12}, \textbf{DI} into \bbhole{pixels.header.data}{J13}, and
\textbf{GND} into \bbhole{pixels.header.ground}{J14}. Then add these connections:
\begin{center}
\begin{tabular}{|l|l|l|}
    \hline
    \textbf{WS2812B pin} & \textbf{XIAO pin} & \textbf{Jumper holes} \\
    \hline
    \connectionrow{pixels-power}{5V (power in)}{3V3}{I3}{I12}
    \connectionrow{pixels-data}{DI (first LED)}{GPIO 20}{I7}{I13}
    \connectionrow{pixels-ground}{GND}{GND}{I2}{I14}
    \hline
\end{tabular}
\end{center}

\subsubsection{Button wiring}
Press the button into the bottom edge of the board so that its legs reach from column \textbf{A}
up to column \textbf{D}. One pair of legs goes into \bbhole{button.side-a}{A12} and \textbf{D12},
and the other pair two rows along into \bbhole{button.side-b}{A14} and \textbf{D14}. The two legs
of a pair are already joined inside the button, so it is the gap between the two pairs that the
button closes. Then add these connections:
\begin{center}
\begin{tabular}{|l|l|l|}
    \hline
    \textbf{Button connection} & \textbf{XIAO pin} & \textbf{Jumper holes} \\
    \hline
    \connectionrow{button-signal}{One side}{GPIO 21}{C7}{E12}
    \connectionrow{button-ground}{Other side}{GND}{J2}{E14}
    \hline
\end{tabular}
\end{center}

The code enables the internal pull-up resistor, so the pin reads HIGH when idle and LOW when pressed.

\begin{figure}[H]
    \centering
    \maybeincludegraphics[width=.9\linewidth]{project_2/wiring.png}
    \caption{Connect the LED strip and button to the XIAO ESP32-C3.}
\end{figure}

\subsection{Programming the microcontroller}
Open the file \texttt{project\_2\_red\_light\_green\_light.py} and run it in the IDE (Chapter \ref{ide}).

At startup, the serial console will print instructions and wait for a button press to begin:
\begin{verbatim}
Project 2: Red Light, Green Light
Press button to start.
\end{verbatim}

During gameplay:
\begin{itemize}
    \item The 3 LEDs act like a stoplight: red LED, yellow LED, and green LED are used individually.
    \item RED phase lights only the red LED and prints ``RED! Do not press.''
    \item YELLOW phase lights only the yellow LED and prints ``YELLOW! Get ready...''
    \item GREEN phase lights only the green LED and prints ``GREEN! Press now.''
    \item Idle/round-over uses the yellow LED as a waiting indicator.
    \item A valid green press prints reaction time and increases score.
    \item A missed green window is counted as one of the 5 windows.
    \item A red/yellow press ends the round and prints final score + best reaction.
\end{itemize}

\subsection{Hardware test checklist}
Use this checklist to confirm your circuit and code are working:
\begin{enumerate}
    \item Confirm LED order on your strip: first LED (DI side) is red, middle is yellow, last is green.
    \item Press button from idle state: yellow wait indicator transitions into red/green gameplay cycling.
    \item Wait for GREEN (green LED only) and press quickly: score increases and a reaction time is printed.
    \item Do not press during RED or YELLOW: pressing before green ends the round immediately.
    \item Skip pressing during one green window: verify it is counted as a miss and the next window starts.
    \item Repeat at least 5 rounds to confirm stable behavior and reliable button detection.
\end{enumerate}

\section{Review}
In this project, you combined timed logic, button input, and LED output to build a full game loop.
You practiced:
\begin{itemize}
    \item State-based programming (idle, green, red, round\_over)
    \item Reusing a shared debounced button library for cleaner input handling
    \item Using \texttt{time.ticks\_ms()} and \texttt{time.ticks\_diff()} for timing
    \item Providing immediate visual and text feedback to players
\end{itemize}

\section{Possible Extensions}
Try these enhancements:
\begin{itemize}
    \item Add difficulty levels by shrinking GREEN timing windows
    \item Display average reaction time after each round
    \item Add a buzzer sound effect for valid and invalid presses
    \item Use different LED patterns (chase, blink, gradient) per game state
    \item Add a countdown before each round starts
\end{itemize}
